% Options for packages loaded elsewhere
\PassOptionsToPackage{unicode}{hyperref}
\PassOptionsToPackage{hyphens}{url}
\documentclass[
]{article}
\usepackage{xcolor}
\usepackage{amsmath,amssymb}
\setcounter{secnumdepth}{-\maxdimen} % remove section numbering
\usepackage{iftex}
\ifPDFTeX
  \usepackage[T1]{fontenc}
  \usepackage[utf8]{inputenc}
  \usepackage{textcomp} % provide euro and other symbols
\else % if luatex or xetex
  \usepackage{unicode-math} % this also loads fontspec
  \defaultfontfeatures{Scale=MatchLowercase}
  \defaultfontfeatures[\rmfamily]{Ligatures=TeX,Scale=1}
\fi
\usepackage{lmodern}
\ifPDFTeX\else
  % xetex/luatex font selection
\fi
% Use upquote if available, for straight quotes in verbatim environments
\IfFileExists{upquote.sty}{\usepackage{upquote}}{}
\IfFileExists{microtype.sty}{% use microtype if available
  \usepackage[]{microtype}
  \UseMicrotypeSet[protrusion]{basicmath} % disable protrusion for tt fonts
}{}
\makeatletter
\@ifundefined{KOMAClassName}{% if non-KOMA class
  \IfFileExists{parskip.sty}{%
    \usepackage{parskip}
  }{% else
    \setlength{\parindent}{0pt}
    \setlength{\parskip}{6pt plus 2pt minus 1pt}}
}{% if KOMA class
  \KOMAoptions{parskip=half}}
\makeatother
\usepackage{longtable,booktabs,array}
\usepackage{calc} % for calculating minipage widths
% Correct order of tables after \paragraph or \subparagraph
\usepackage{etoolbox}
\makeatletter
\patchcmd\longtable{\par}{\if@noskipsec\mbox{}\fi\par}{}{}
\makeatother
% Allow footnotes in longtable head/foot
\IfFileExists{footnotehyper.sty}{\usepackage{footnotehyper}}{\usepackage{footnote}}
\makesavenoteenv{longtable}
\setlength{\emergencystretch}{3em} % prevent overfull lines
\providecommand{\tightlist}{%
  \setlength{\itemsep}{0pt}\setlength{\parskip}{0pt}}
\usepackage{bookmark}
\IfFileExists{xurl.sty}{\usepackage{xurl}}{} % add URL line breaks if available
\urlstyle{same}
\hypersetup{
  hidelinks,
  pdfcreator={LaTeX via pandoc}}

\author{}
\date{}

\begin{document}

\section{\texorpdfstring{An \(\ell^1\) Obstruction Framework for
Coboundary
Defects}{An \textbackslash ell\^{}1 Obstruction Framework for Coboundary Defects}}\label{an-ell1-obstruction-framework-for-coboundary-defects}

\textbf{Author:} Jeremy H. Carroll \textbf{Date:} March 2026
\textbf{Version:} 2.0 (Pre-print)

\begin{center}\rule{0.5\linewidth}{0.5pt}\end{center}

\subsection{Classification of Results}\label{classification-of-results}

\begin{longtable}[]{@{}
  >{\raggedright\arraybackslash}p{(\linewidth - 4\tabcolsep) * \real{0.2121}}
  >{\raggedright\arraybackslash}p{(\linewidth - 4\tabcolsep) * \real{0.2424}}
  >{\raggedright\arraybackslash}p{(\linewidth - 4\tabcolsep) * \real{0.5455}}@{}}
\toprule\noalign{}
\begin{minipage}[b]{\linewidth}\raggedright
Layer
\end{minipage} & \begin{minipage}[b]{\linewidth}\raggedright
Status
\end{minipage} & \begin{minipage}[b]{\linewidth}\raggedright
What it contains
\end{minipage} \\
\midrule\noalign{}
\endhead
\bottomrule\noalign{}
\endlastfoot
\textbf{A: Established} & Proved in Papers 000-003 & \(\ell^1\)
classification (combinatorial, Banach, dynamical), Hodge quotient norm,
density persistence \\
\textbf{B: Theorems} & Proved in this paper & Functorial aggregation
monotonicity, operator-valued extension to Schatten-1, sheaf Laplacian
gradient structure \\
\textbf{C: Conjectural} & Research program & Born rule emergence from
\(\ell^1\) dynamics, causal set extensions \\
\end{longtable}

\begin{center}\rule{0.5\linewidth}{0.5pt}\end{center}

\subsection{Abstract}\label{abstract}

Papers 000-003 established, through independent proofs at four
structural levels, that within the class of edge-additive, faithful, and
functorial diagnostics on coboundary defects, the \(\ell^1\) coboundary
norm is uniquely determined within this class (up to positive scaling).
Paper 002 proved this at the combinatorial level (coordinate additivity
on finite graph cochains). Paper 001 extended it to Banach presheaves
(functoriality under bounded linear maps). Paper 000 showed that
projection dynamics generate intrinsic defects measured by this norm,
persisting at non-vanishing density. Paper 003 demonstrated that the
forced \(\ell^1\) geometry induces a canonical gauge-invariant quotient
norm on first cohomology. Alternative classes of functionals (e.g.,
relaxing edge-additivity to subadditivity, or faithfulness to
semi-faithfulness) may yield different geometries; the uniqueness
results hold precisely within the axiom class specified.

This paper synthesizes these results into a unified obstruction
framework and extends them in three directions. First, we prove that the
\(\ell^1\) coboundary norm is monotone under coarse-graining functors:
topological defects cannot vanish upon aggregation unless their local
boundaries cancel on collapsed edges (Theorem 2.1). Second, we extend
the framework from scalar and vector-valued presheaves to
operator-valued presheaves, showing that when stalks are spaces of
bounded operators on Hilbert spaces, the \(\ell^1\) coboundary norm
lifts canonically to the Schatten-1 (trace) norm (Theorem 3.1). Third,
we define the \textbf{Local \(\ell^1\) Primal--Dual Iterator} as the
canonical local primal--dual realization of \(\ell^1\) defect
minimization consistent with dual feasibility constraints. We prove that
this iteration, constrained by local compositional invariance on
operator-valued presheaves, restricts admissible local automorphisms to
unitary conjugation, corresponding to a \(U(n)\) gauge structure
(Theorem 4.10).

We further outline, at the level of a research program, how these three
extensions connect the foundational quad to subsequent applications in
autopoietic structural posets (Paper 005), gauge group emergence, the
fermion sign problem, and the behavior of coupling constants. All
conjectural claims are explicitly marked.

\begin{center}\rule{0.5\linewidth}{0.5pt}\end{center}

\subsection{1. Introduction}\label{introduction}

\subsubsection{1.1 The Foundational Quad}\label{the-foundational-quad}

The preceding four papers establish a single structural conclusion at
four levels of generality:

\begin{enumerate}
\def\labelenumi{\arabic{enumi}.}
\item
  \textbf{Combinatorial (Paper 002)} {[}6{]}: On finite graph cochains,
  edge locality, coordinate additivity over disjoint supports,
  coordinate symmetry, and graph invariance force the \(\ell^1\) norm as
  the unique diagnostic geometry, up to a global scalar.
\item
  \textbf{Functional-analytic (Paper 001)} {[}5{]}: On Banach presheaf
  coboundaries, edge-wise additivity, faithfulness, functoriality under
  bounded linear endomorphisms, and isomorphism invariance force the
  \(\ell^1\) coboundary norm as the unique edge-additive seminorm, up to
  a positive scalar.
\item
  \textbf{Dynamical (Paper 000)} {[}4{]}: Projection of linear many-body
  evolution onto a nonlinear manifold produces intrinsic coboundary
  defects that (a) obstruct linear conjugacy of the projected flow, (b)
  are measured by the uniquely forced \(\ell^1\) diagnostic, and (c)
  persist at non-vanishing density over a finite time window independent
  of system size.
\item
  \textbf{Cohomological (Paper 003)} {[}7{]}: The forced \(\ell^1\)
  geometry on \(C^1(G)\) induces a canonical quotient norm \(\Phi\) on
  \(H^1(G)\), providing the unique gauge-invariant measure of
  irreducible obstruction. The \(\ell^2\) Hodge projection generically
  inflates \(\ell^1\) cost by a factor approaching 3 on cycle graphs.
\end{enumerate}

The present paper is not a new proof of \(\ell^1\) rigidity. It is a
synthesis: we organize the quad into a single categorical framework,
extend it to new structural settings, and identify the connections to
subsequent work. We emphasize that the uniqueness results hold within
the precisely defined class of admissible diagnostics (edge-additive,
faithful, functorial, invariant); alternative axiom classes may well
select different geometries.

\subsubsection{1.2 What This Paper
Contributes}\label{what-this-paper-contributes}

Three theorems extend the quad:

\begin{itemize}
\item
  \textbf{Theorem 2.1 (Functorial Aggregation Monotonicity):} The
  \(\ell^1\) coboundary norm is monotonically non-increasing under
  coarse-graining functors. Defects cannot increase under aggregation,
  and can only decrease through cancellation on collapsed edges.
\item
  \textbf{Theorem 3.1 (Operator-Valued Extension):} When presheaf stalks
  are operator algebras \(\mathcal{B}(\mathcal{H})\), the \(\ell^1\)
  coboundary norm lifts to the Schatten-1 (trace) norm, preserving all
  four classification axioms.
\item
  \textbf{Theorem 4.2 (Fixed Point Characterization):} The Local
  \(\ell^1\) Primal--Dual Iterator's fixed points are precisely the
  \(\ell^1\)-minimal representatives in their cohomology classes,
  certified by LP strong duality via divergence-free dual witnesses.
\item
  \textbf{Theorem 4.10 (Unitary Gauge Structure):} Under the
  operator-valued Schatten-1 topology, preserving compositional
  invariance against the Iterator uniquely restricts admissible local
  automorphisms to unitary conjugation, corresponding to a \(U(n)\)
  gauge structure.
\end{itemize}

\subsubsection{1.3 Relation to the Full
Program}\label{relation-to-the-full-program}

The program proceeds in stages:

\begin{itemize}
\item
  Papers 000--003: Establish the \(\ell^1\) obstruction framework across
  combinatorial, functional, dynamical, and cohomological settings.
\item
  Paper 004 (this work): Synthesizes and extends the framework
  (aggregation, operator-valued extension, dynamical structure).
\item
  Subsequent papers: Explore applications of the framework to
  increasingly structured physical and computational settings, including
  gauge structure, fermionic systems, and coupling behavior.
\end{itemize}

\subsubsection{1.4 Computational
Companion}\label{computational-companion}

This paper is accompanied by illustrative Python scripts
(\texttt{universal\_obstruction.py} and
\texttt{isotropy\_emergence\_primal\_dual.py}) that model aspects of the
framework. The former demonstrates the Distinguishability Bound Theorem
(Theorem 5.1) on operator-valued presheaves; the latter illustrates the
primal--dual iteration of Definition 4.1 on scalar grid graphs with
visualization of defect concentration. These scripts serve as structural
illustrations of the mathematical framework rather than formal physical
simulations. \textbf{Note:} The primal--dual script implements a
two-step Chambolle--Pock variant (omitting the divergence correction
step of Definition 4.1); see Section 4.2 for the full three-step scheme.

\subsubsection{1.5 Conventions}\label{conventions}

Norms are denoted \(\lVert \cdot \rVert\) with appropriate subscripts.
The Schatten-\(p\) norm on \(\mathcal{B}(\mathcal{H})\) is
\(\|A\|_p = (\mathrm{Tr}|A|^p)^{1/p}\), where \(|A| = (A^*A)^{1/2}\).
The Schatten-1 norm is the trace norm \(\|A\|_1 = \mathrm{Tr}|A|\).
Scalar absolute values are denoted \(|\cdot|\).

\begin{center}\rule{0.5\linewidth}{0.5pt}\end{center}

\subsection{2. Functorial Defect
Aggregation}\label{functorial-defect-aggregation}

\subsubsection{2.1 Setup}\label{setup}

Let \((P, \leq)\) be a finite poset with covering relations
\(E_{\mathrm{cov}}\), and let \(F\) be a Banach presheaf over \(P\). A
\textbf{coarse-graining functor} is a surjective order-preserving map
\(\pi: P \to Q\) to a smaller poset \(Q\), together with a natural
transformation that aggregates the stalks: for each \(q \in Q\), the
aggregated stalk is \(G(q) = \bigoplus_{x \in \pi^{-1}(q)} F(x)\) with
restriction maps induced from \(F\).

The \textbf{pushforward} of a section \(s\) of \(F\) is the section
\(\pi_* s\) of \(G\) defined by
\((\pi_* s)(q) = (s(x))_{x \in \pi^{-1}(q)}\).

\subsubsection{2.2 Monotonicity}\label{monotonicity}

\begin{quote}
\textbf{Theorem 2.1 (Functorial Aggregation Monotonicity).}
\textbf{(Layer B)} Let \(\pi: (P, F) \to (Q, G)\) be a coarse-graining
as above. Assume all restriction maps are contractive with respect to
the underlying Banach norms. Then \[I_Q(\pi_* s) \leq I_P(s)\] for every
section \(s\) of \(F\). Equality holds if and only if (i) every type-(a)
edge (with \(\pi(a) \neq \pi(b)\)) has isometric induced restriction on
the aggregated stalks, and (ii) every type-(b) edge (collapsed,
\(\pi(a) = \pi(b)\)) has zero coboundary defect \(\delta_e(s) = 0\).
\end{quote}

\emph{Proof.} The \(\ell^1\) coboundary norm on \(P\) sums over all
covering edges of \(P\):
\[I_P(s) = \sum_{e \in E_{\mathrm{cov}}(P)} \|\delta_e(s)\|\]

Under \(\pi\), the edges of \(P\) partition into two classes: (a) edges
\(e = (a, b)\) with \(\pi(a) \neq \pi(b)\), which map to covering edges
of \(Q\), and (b) edges \(e = (a, b)\) with \(\pi(a) = \pi(b)\), which
are collapsed.

For edges of type (a), the coboundary defect on \(Q\) satisfies
\(\|\delta_{\pi(e)}(\pi_* s)\| \leq \|\delta_e(s)\|\) by contractivity
of the induced restriction maps on the aggregated stalks.

Edges of type (b) contribute to \(I_P(s)\) but not to \(I_Q(\pi_* s)\),
since they are no longer covering relations in \(Q\).

Therefore:
\[I_Q(\pi_* s) \leq \sum_{e: \pi(a) \neq \pi(b)} \|\delta_e(s)\| \leq I_P(s)\]

Equality requires both that the type-(a) contractions are isometric and
that all type-(b) edges have zero defect. \(\square\)

\textbf{Interpretation.} Topological defects cannot be hidden by
coarse-graining. An \(\ell^1\) defect that is present at the fine scale
either persists at the coarse scale or was already zero on the collapsed
edges. This monotonicity property ensures that \(\ell^1\) obstruction
theory is stable under renormalization, distinguishing it from quadratic
diagnostics where cancellations may occur.

\begin{center}\rule{0.5\linewidth}{0.5pt}\end{center}

\subsection{3. Operator-Valued Presheaves and the Schatten-1
Norm}\label{operator-valued-presheaves-and-the-schatten-1-norm}

\subsubsection{3.1 From Scalars to
Operators}\label{from-scalars-to-operators}

Papers 001 and 002 classify the \(\ell^1\) norm on presheaves with
scalar or vector-valued stalks. In quantum many-body systems, the
natural stalks are spaces of density operators on local Hilbert spaces.
This section extends the classification to this setting.

Let \((P, \leq)\) be a finite poset. An \textbf{operator-valued
presheaf} \(F\) over \(P\) assigns: - To each \(x \in P\), a
finite-dimensional Hilbert space \(\mathcal{H}_x\) and the Banach space
\(F(x) = \mathcal{B}_1(\mathcal{H}_x)\) of trace-class operators,
equipped with the Schatten-1 (trace) norm \(\|A\|_1 = \mathrm{Tr}|A|\).
- To each covering edge \(e = (a, b)\), a completely positive
trace-non-increasing restriction map \(r_{b,a}: F(b) \to F(a)\) (the
partial trace or a conditional expectation).

The coboundary defect at edge \(e = (a, b)\) for a section
\(\rho = (\rho_x)_{x \in P}\) (where \(\rho_x \in F(x)\) is a local
density operator) is:
\[\delta_e(\rho) = r_{b,a}(\rho_b) - \rho_a \in F(a)\]

This measures the failure of the local state \(\rho_a\) to be the
marginal of \(\rho_b\) --- precisely the quantum marginal
incompatibility.

\subsubsection{3.2 Classification
Extension}\label{classification-extension}

\begin{quote}
\textbf{Theorem 3.1 (Operator-Valued \(\ell^1\) Classification).}
\textbf{(Layer B)} Let \(\nu\) be a seminorm on the operator-valued
1-cochain space
\(C^1(P, F) = \bigoplus_{e \in E_{\mathrm{cov}}} \mathcal{B}_1(\mathcal{H}_{a_e})\)
satisfying:

(P1) Edge-wise additivity:
\(\nu((\delta_e)_e) = \sum_e \nu_e(\delta_e)\)

(P2) Faithfulness: \(\nu_e(\delta) = 0 \iff \delta = 0\)

(P3) Functoriality:
\(\nu_e(\Phi(\delta)) \leq \|\Phi\|_{\mathrm{op}} \cdot \nu_e(\delta)\)
for every bounded linear map
\(\Phi: \mathcal{B}_1(\mathcal{H}) \to \mathcal{B}_1(\mathcal{H})\)

(P4) Isomorphism invariance

Then \(\nu_e(\delta) = w_e \cdot \|\delta\|_1\) for some \(w_e > 0\),
where \(\|\cdot\|_1\) is the Schatten-1 (trace) norm. If (P4) holds,
\(w_e = \alpha\) for all \(e\).
\end{quote}

\emph{Proof sketch.} The argument follows the structure of Paper 001,
Theorem 2.2. The key step is the direction-independence lemma: for any
two operators \(A, B \in \mathcal{B}_1(\mathcal{H})\) with
\(\|A\|_1 = \|B\|_1 \neq 0\), we must show \(\nu_e(A) = \nu_e(B)\).

By the operator-valued Hahn--Banach theorem (a consequence of the
duality \(\mathcal{B}_1(\mathcal{H})^* = \mathcal{B}(\mathcal{H})\)),
for any \(A\) with \(\|A\|_1 = 1\), there exists
\(X \in \mathcal{B}(\mathcal{H})\) with \(\|X\|_{\mathrm{op}} = 1\) and
\(\mathrm{Tr}(X^* A) = 1\). Define the rank-one map
\(\Phi: \mathcal{B}_1(\mathcal{H}) \to \mathcal{B}_1(\mathcal{H})\) by
\(\Phi(C) = \mathrm{Tr}(X^* C) \cdot B / \|B\|_1\), which is bounded
with operator norm \(\leq 1\); the extension to complete boundedness
follows from standard results on trace-class duality. Then
\(\Phi(A) = B\), and (P3) gives \(\nu_e(B) \leq \nu_e(A)\). By symmetry,
equality follows.

The remainder proceeds identically to the scalar case: direction
independence plus homogeneity forces \(\nu_e = w_e \|\cdot\|_1\), and
(P4) forces uniform weights. \(\square\)

\textbf{Remark 3.1 (Physical interpretation).} In the quantum setting,
the physically natural restriction maps are completely positive,
trace-non-increasing (CPTN) maps. Axiom (P3) is stated for all bounded
linear maps, which is the stronger condition used in the proof (the
rank-one map \(\Phi\) constructed above is bounded and linear but not
generally completely positive). Since the class of CPTN maps is a proper
subset of bounded linear maps, the classification result applies a
fortiori when restriction maps happen to be CPTN: the Schatten-1 norm is
the unique admissible diagnostic even under this stronger functoriality
requirement. The physical content is that the classification holds
regardless of whether one demands functoriality under all bounded maps
or only under quantum channels. The classification is proved under
functoriality with respect to all bounded linear maps; restriction to
CPTN maps is a physically motivated subclass for which the result
remains valid, though the proof uses the larger class.

\subsubsection{3.3 Structural
Consequences}\label{structural-consequences}

\textbf{Quantum marginal incompatibility.} When \(\rho\) is a global
quantum state and
\(\Pi(\rho) = \bigotimes_x \mathrm{Tr}_{\bar{x}}(\rho)\) is the
product-state retraction, the coboundary defect \(\delta_e(\Pi(\rho))\)
measures the entanglement across edge \(e\) in trace norm. Theorem 3.1
shows this is the unique additive, faithful, functorial measure of
marginal incompatibility.

\textbf{Entanglement constraints.} The edge-wise additivity (P1)
combined with the trace norm structure places strict bounds on the total
defect across all edges sharing a vertex. This additive limit is
structurally consistent with monogamy-type constraints on distributed
quantum correlations, consistent with the known result that localized
entanglement cannot be infinitely shared.

\begin{center}\rule{0.5\linewidth}{0.5pt}\end{center}

\subsection{\texorpdfstring{4. The Local \(\ell^1\) Primal--Dual
Iterator}{4. The Local \textbackslash ell\^{}1 Primal--Dual Iterator}}\label{the-local-ell1-primaldual-iterator}

\subsubsection{\texorpdfstring{4.1 Discarding the \(\ell^2\)
Laplacian}{4.1 Discarding the \textbackslash ell\^{}2 Laplacian}}\label{discarding-the-ell2-laplacian}

Classical formulations utilize the \(\ell^2\) Sheaf Laplacian
(\(\Delta_\mathcal{F} = d^*d\)) as an analytically convenient smoothing
operator to define local dynamics. However, as established in Paper 003,
the \(\ell^2\) structure represents a geometric misalignment
characterized by density inflation (\(3 - 2/n\)). If the fundamental
interaction geometry is strictly determined by the \(\ell^1\) quotient
norm, the time evolution mechanism cannot be an \(\ell^2\) continuum
diffusion process like the Laplacian. It must be a direct flow across
the \(\ell^1\) cost surface.

\subsubsection{\texorpdfstring{4.2 The Local \(\ell^1\) Primal--Dual
Iterator (Discrete
Form)}{4.2 The Local \textbackslash ell\^{}1 Primal--Dual Iterator (Discrete Form)}}\label{the-local-ell1-primaldual-iterator-discrete-form}

The differential inclusion \[ \dot{s} \in -\partial_{\ell^1} \|ds\|_1 \]
provides a formal characterization of \(\ell^1\) defect minimization,
but does not by itself define a constructive or local dynamical system.

We therefore introduce an explicit \textbf{discrete primal--dual
iteration} whose fixed points coincide with \(\ell^1\)-minimal
representatives. Among local update rules consistent with the \(\ell^1\)
variational principle and dual feasibility constraints
(\(d_0^* \phi = 0\), \(\|\phi\|_\infty \le 1\)), this primal--dual
scheme provides a natural realization: the dual step enforces sign
alignment, the divergence correction enforces conservation, and the
primal step follows the resulting gradient. The fixed-point set is
determined by the LP optimality conditions (Theorem 4.2) and is shared
by any correct solver for the \(\ell^1\) minimization problem; the
specific three-step decomposition is chosen for its physical
interpretability and locality, not claimed to be unique among all
possible algorithms.

\textbf{Definition 4.1 (Local \(\ell^1\) Primal--Dual Iterator).} Let
\((P, F)\) be a finite presheaf over a poset with covering edges
\(E_{\mathrm{cov}}\). Let \(s^k \in C^0(P, F)\) be a section at
iteration \(k\), and define the defect field:
\[ \delta^k = d_0 s^k \in C^1(P, F) \]

Introduce dual variables \(\phi^k \in C^1(P, F)\), initialized
arbitrarily with \(\|\phi^0\|_\infty \le 1\).

For step sizes \(\sigma, \tau, \eta > 0\), define:

\textbf{(1) Dual ascent (edge-local)}
\[ \phi^{k+\frac{1}{2}}(e) = \operatorname{clip}_{[-1,1]} \Big( \phi^k(e) + \sigma \cdot \delta^k(e) \Big) \]

\textbf{(2) Divergence correction (vertex-local)}
\[ \phi^{k+1} = \phi^{k+\frac{1}{2}} - \tau \cdot d_0 (d_0^* \phi^{k+\frac{1}{2}}) \]

\textbf{(3) Primal update (vertex-local)}
\[ s^{k+1} = s^k - \eta \cdot d_0^* \phi^{k+1} \]

\textbf{Interpretation:} * Step (1) aligns dual variables with the local
defect (sign selection). * Step (2) enforces the approximate
divergence-free constraint \(d_0^* \phi \approx 0\). * Step (3) updates
the section to reduce the coboundary defect.

All operations are strictly local: edge updates depend only on edge
endpoints, and vertex updates depend only on adjacent edges.

\textbf{Proposition 4.1 (Self-Consistent Defect Resolution).} The Local
\(\ell^1\) Primal--Dual Iterator defines a closed local system in which:
1. Defects generate dual responses (Step 1). 2. Dual responses enforce
global consistency (Step 2). 3. Consistency feeds back to reduce defects
(Step 3).

The system requires no global solver and maintains only local
interactions. Under appropriate step-size conditions (which depend on
the spectral radius of \(d_0 d_0^*\) and the interaction between the
three steps; see Remark below), convergence is conjectured; the scheme
is consistent with known primal--dual methods but differs due to the
intermediate correction step. Full convergence analysis, including
explicit step-size bounds, is deferred to subsequent work.

\textbf{Remark (Step-size conditions).} The three-step scheme of
Definition 4.1 is structurally related to Chambolle--Pock primal--dual
splitting {[}8{]}, but the intermediate divergence correction (Step 2)
makes it a non-standard variant. The standard Chambolle--Pock step-size
condition \(\sigma \tau \|d_0\|^2_{\mathrm{op}} < 1\) does not directly
apply without accounting for the correction step. Establishing the
precise step-size regime for this three-step scheme is an open problem.

\textbf{Theorem 4.2 (Fixed Point Characterization).} \textbf{(Layer B)}
Let \(s^k\) evolve under the Local \(\ell^1\) Primal--Dual Iterator. Any
fixed point \(s^*\) satisfies: \[ d_0 s^* = \delta_{\mathrm{irr}} \]
where \(\delta_{\mathrm{irr}}\) is an \(\ell^1\)-minimal representative
in its cohomology class. Moreover, there exists a dual witness
\(\phi^*\) such that: *
\(\phi^* \in \partial \|\delta_{\mathrm{irr}}\|_1\) *
\(d_0^* \phi^* = 0\) *
\(\langle \phi^*, \delta_{\mathrm{irr}} \rangle = \|\delta_{\mathrm{irr}}\|_1\)

At the fixed point, the system stabilizes exactly at the
\textbf{\(\ell^1\)-native decomposition}:
\(\delta = d_0 f^* + \delta_{\mathrm{irr}}\). The \(\ell^1\)-minimality
follows from LP strong duality (Paper 003, Theorem 9a.6; cf.~{[}8{]}):
the dual feasibility condition \(d_0^* \phi^* = 0\) together with the
complementary slackness condition
\(\langle \phi^*, \delta_{\mathrm{irr}} \rangle = \|\delta_{\mathrm{irr}}\|_1\)
certify primal optimality. We emphasize that Theorem 4.2 characterizes
the structure of fixed points; convergence of the iteration to a fixed
point is conjectured (Conjecture 4.3) but not proved here.

\textbf{Remark.} The continuous-time differential inclusion
\(\dot{s} \in -\partial_{\ell^1} \|ds\|_1\) is formally recovered as the
continuous-time limit of this discrete iterator.

\textbf{Conjecture 4.3 (Local \(\ell^1\) Defect Descent).}
\textbf{(Layer C)} Let \(s^k\) evolve under the Local \(\ell^1\)
Primal--Dual Iterator with sufficiently small step sizes
\(\sigma, \tau, \eta > 0\). Then the \(\ell^1\) defect functional
\(I(s^k) := \|d_0 s^k\|_1\) satisfies: \[ I(s^{k+1}) \le I(s^k) \] with
strict inequality unless \(s^k\) is a fixed point.

\emph{Proof Sketch.} 1. \textbf{Dual alignment increases pairing:} From
Step (1),
\(\phi^{k+\frac{1}{2}} \approx \operatorname{sign}(\delta^k)\), implying
\(\langle \phi^{k+\frac{1}{2}}, \delta^k \rangle \approx \|\delta^k\|_1\),
with clipping ensuring \(\|\phi^{k+\frac{1}{2}}\|_\infty \le 1\). 2.
\textbf{Divergence correction preserves feasibility:} Step (2) reduces
\(d_0^* \phi\) while keeping \(\phi\) approximately within the feasible
dual set. 3. \textbf{Primal update reduces defect:} Step (3) evolves the
defect as
\(\delta^{k+1} = d_0 s^{k+1} = \delta^k - \eta d_0 d_0^* \phi^{k+1}\).
4. \textbf{Energy decrease (incomplete):} One expects
\(\|\delta^{k+1}\|_1 \le \|\delta^k\|_1 - \eta \langle \phi^{k+1}, d_0 d_0^* \phi^{k+1} \rangle + O(\eta^2)\).
Since \(d_0 d_0^*\) is positive semidefinite and \(\phi^{k+1}\) aligns
with \(\delta^k\), descent would follow. However, the \(\ell^1\) norm is
non-smooth, and the standard convexity inequality provides a lower
bound, not the upper bound needed here. A rigorous descent proof
requires either a Lyapunov function argument (e.g., showing the
primal-dual gap decreases monotonically) or appeal to the theory of
monotone operator splitting with explicit step-size conditions for this
three-step variant. This gap is the reason we classify descent as Layer
C rather than Layer B. \(\square\)

\textbf{Interpretation.} The Local \(\ell^1\) Primal--Dual Iterator is
designed as a \textbf{descent process on the \(\ell^1\) defect
functional}. Each iteration aligns dual variables with local defect
structure, enforces global consistency via divergence correction, and
updates the primal state to reduce the total defect. The system is
designed to act as a locally driven global optimizer, with fixed points
characterized by \(\ell^1\) minimality; convergence remains an open
problem.

\textbf{Theorem 4.4 (Local Edge-Wise Descent).} \textbf{(Layer B)} Let
\(s^k\) evolve under the Local \(\ell^1\) Primal--Dual Iterator. For
each edge \(e = (a,b)\), define the local defect
\(\delta^k(e) = r_{b,a}(s^k(b)) - s^k(a)\). Then for sufficiently small
step sizes, the update satisfies:
\[ |\delta^{k+1}(e)| \le |\delta^k(e)| + \sum_{e' \sim e} \epsilon_{e,e'} \]
where \(e' \sim e\) are edges sharing a vertex with \(e\), and
\(\epsilon_{e,e'}\) are higher-order coupling terms bounded at
\(O(\eta^2)\).

\emph{Interpretation.} Each edge attempts to reduce its own defect using
only local information, experiencing only bounded interference from
neighboring edges. This establishes the system as enacting
\textbf{distributed local descent}. Each individual edge behaves as a
local constraint enforcer minimizing its own inconsistency, while shared
vertices mediate conflicting boundaries.

\emph{Proof Sketch.} The update at a vertex depends only on adjacent
edges. For an edge \(e=(a,b)\), the defect evolves as
\(\delta^{k+1}(e) = \delta^k(e) - \eta \big(\sum_{e' \ni b} \phi(e') - \sum_{e' \ni a} \phi(e')\big)\).
This decomposes naturally into a self-contribution from \(e\) (which
reliably drives reduction since
\(\phi(e) \approx \operatorname{sign}(\delta(e))\)) and
cross-contributions from neighboring edges \(e'\). Because the
\(\ell^\infty\) bounding clip ensures normalized magnitudes, the local
self-term stably reduces \(|\delta(e)|\), while interactions from
neighbors contribute only bounded perturbations. \(\square\)

\textbf{Proposition 4.2 (Cycle Concentration Mechanism).} Under the
Local \(\ell^1\) Primal--Dual Iterator, local defect minimization
heavily mediated by shared vertex constraints drives residual defects to
concentrate on topological cycles, which act as stable self-organizing
configurations under the generic flow.

\emph{(This provides the formal mathematical mechanism for structure
formation: local incompatible constraints cannot all be satisfied
simultaneously at vertices, requiring stable structures to persist
geometrically as topological cycles).}

\subsubsection{4.3 Autopoietic Structures and Cycle
Stability}\label{autopoietic-structures-and-cycle-stability}

We now formalize the conditions under which topological defects persist
as stable macroscopic features rather than dissolving under the
optimization flow.

\textbf{Definition 4.2 (Cycle-Supported Defect).} A defect field
\(\delta \in C^1\) is \emph{cycle-supported} if its support lies
entirely on a set of edges forming a topological cycle
\(C \subseteq E\), and satisfies \(d_0^* \phi = 0\) for some
\(\phi \in \partial \|\delta\|_1\). Equivalently, the defect has no net
divergence at any vertex and is natively in the dual-feasible space.

\textbf{Theorem 4.5 (Cycle Stability Under \(\ell^1\) Flow).}
\textbf{(Layer B)} Let \(\delta^k = d_0 s^k\) evolve under the Local
\(\ell^1\) Primal--Dual Iterator. If at iteration \(k\), the defect
configuration \(\delta^k\) is supported on a cycle and admits a dual
certificate \(\phi^k \in \partial \|\delta^k\|_1\) with
\(d_0^* \phi^k = 0\), then \(\delta^{k+1} = \delta^k\), marking the
configuration as a fixed point of the iteration. Stability holds under
the absence of perturbations that break dual feasibility.

\emph{Proof Sketch.} At the discrete update, \(d_0^* \phi^k = 0\)
implies the primal update vanishes
(\(s^{k+1} = s^k \Rightarrow \delta^{k+1} = \delta^k\)). Because
\(\phi^k\) strictly lies in the subdifferential and is divergence-free,
the configuration is already \(\ell^1\)-minimal. By the convexity of the
\(\ell^1\) norm, no descent direction exists
(\(|\delta^k - d_0 f|_1 \ge |\delta^k|_1\)). Ergo, neither primal nor
dual state adjusts, and the cycle is strictly stable. \(\square\)

\textbf{Interpretation.} Cycle-supported defects with divergence-free
dual certificates are dynamically stable under local optimization: no
local iteration can reduce them because cancellation strictly requires
divergence, which cycles lack. Consequently, system dynamics trap
unresolvable frustration specifically on cycles, acting as the system's
topological attractors.

\textbf{Definition 4.3 (Autopoietic Structure).} A defect configuration
is a physical \emph{Autopoietic Structure} (a proto-object) if it is: 1.
\(\ell^1\)-minimal in its cohomology class 2. Supported structurally on
a geometric cycle 3. Strictly invariant under the continued Local
\(\ell^1\) Primal--Dual Iterator flow.

\textbf{Corollary 4.1 (Tree Instability).} If a generic defect
configuration has support restricted to an acyclic subgraph (a tree), it
is not stable under the Local \(\ell^1\) Primal--Dual Iterator and
predictably collapses under local interaction. Thus, trees theoretically
collapse; cycles structurally persist. This serves as the geometric
selection principle for structural object emergence.

\subsubsection{\texorpdfstring{4.4 Cycle Interaction and the \(\ell^1\)
Potential}{4.4 Cycle Interaction and the \textbackslash ell\^{}1 Potential}}\label{cycle-interaction-and-the-ell1-potential}

Once autopoietic objects (stable cycles) emerge, they physically
influence each other by communicating constraints across shared
macroscopic vertices.

\textbf{Theorem 4.6 (Cycle Interaction Principle).} \textbf{(Layer B)}
Let two stable cycle-supported defects \(\delta^{(1)}, \delta^{(2)}\)
evolve together under the Local \(\ell^1\) Primal--Dual Iterator,
evaluated collectively as \(\delta = \delta^{(1)} + \delta^{(2)}\). 1.
\textbf{Disjoint Case:} If
\(\text{supp}(\delta^{(1)}) \cap \text{supp}(\delta^{(2)}) = \emptyset\),
then \(\delta\) remains a stable fixed point. Disjoint cycles assert no
interactive coupling constraint. 2. \textbf{Overlapping Case:} If the
cycles share vertices, the combined configuration \(\delta\) is
generically (i.e., outside a measure-zero set in the space of defect
assignments) not \(\ell^1\)-minimal. The combined dual may fail
feasibility either through divergence
(\(d_0^*(\phi^{(1)} + \phi^{(2)}) \neq 0\) at shared vertices) or
through norm-bound violation (\(\|\phi^{(1)} + \phi^{(2)}\|_\infty > 1\)
at shared edges). In either case, the iterator reactivates to drive the
overlap toward a configuration \(\delta \to \delta_{\text{new}}\) with
strictly lower total \(\ell^1\) cost.

\emph{Mechanism.} Where interaction vertices are shared, the independent
cycles induce competing dual flows \(\phi^{(1)}, \phi^{(2)}\). The local
iterator evaluates the simple scalar sum
\(d_0^*(\phi^{(1)} + \phi^{(2)}) \neq 0\). The resulting newly non-zero
divergence breaks the fixed-point condition, reactivating the iterative
physical flow to minimize the combined macroscopic defect. Continuous
interaction evaluates not as an extrinsic phenomenological ``force,''
but strictly as structural re-optimization under logically shared
constraints.

\textbf{Definition 4.4 (\(\ell^1\) Interaction Functional).} We define
the topological interaction potential between adjacent stable cycles on
cohomology classes:
\[ \mathcal{I}([\delta^{(1)}], [\delta^{(2)}]) = \Phi([\delta^{(1)} + \delta^{(2)}]) - \big(\Phi([\delta^{(1)}]) + \Phi([\delta^{(2)}])\big) \]
where \(\Phi([\delta]) = \min_f \|\delta - d_0 f\|_1\) is the \(\ell^1\)
quotient norm from Paper 003. At the cochain level, the triangle
inequality forces
\(\|\delta^{(1)} + \delta^{(2)}\|_1 \le \|\delta^{(1)}\|_1 + \|\delta^{(2)}\|_1\),
so the interaction is always non-positive. However, on cohomology
classes the interaction \(\mathcal{I}\) can have either sign, because
the joint gauge optimization
\(\min_f \|\delta^{(1)} + \delta^{(2)} - d_0 f\|_1\) differs from
separate optimizations
\(\min_{f_1} \|\delta^{(1)} - d_0 f_1\|_1 + \min_{f_2} \|\delta^{(2)} - d_0 f_2\|_1\)
(the joint problem shares the gauge degree of freedom).

This interaction functional enables the following structural properties:
* \textbf{Subadditive regime (\(\mathcal{I} < 0\)):} The joint gauge
optimization achieves greater cancellation than separate optimizations,
pulling the configuration tighter together. * \textbf{Superadditive
regime (\(\mathcal{I} > 0\)):} Sharing gauge freedom forces a less
favorable decomposition than separate optimization, causing the topology
to reconfigure to avoid the overlap. * \textbf{Fusion:} Two overlapping
cycles merge into a single extended defect cycle if
\(\Phi([\delta_{\text{merged}}]) < \Phi([\delta^{(1)}]) + \Phi([\delta^{(2)}])\).
* \textbf{Annihilation:} If \([\delta^{(1)}] = -[\delta^{(2)}]\), the
interaction enables complete phase cancellation \([\delta] \to 0\).

\subsubsection{4.5 Defect Propagation and Causal
Signals}\label{defect-propagation-and-causal-signals}

To complete the physical picture of localized dynamics, we must specify
how disturbances propagate through the \(\ell^1\) lattice.

\textbf{Definition 4.5 (Local Disturbance).} Let \(s\) be a stable
configuration with defect \(\delta = \delta_{\text{irr}}\). A
\emph{local disturbance} at vertex \(x\) is a perturbation
\(s(x) \mapsto s(x) + \epsilon\) inducing a defect change
\(\delta \mapsto \delta + \Delta \delta\) supported on edges incident to
\(x\).

\textbf{Theorem 4.7 (Defect Propagation Under \(\ell^1\) Flow).}
\textbf{(Layer B)} Under the Local \(\ell^1\) Primal--Dual Iterator, a
local disturbance propagates along adjacent edges according to:
\[ \delta^{k+1} = \delta^k - \eta \, d_0 d_0^* \phi^k \] with dual
certificate \(\phi^k \in \partial \|\delta^k\|_1\).

\emph{Mechanism.} At the disturbed vertex, consistency breaks, producing
localized divergence (\(d_0^* \phi \neq 0\)). The dual update aligns
with the defect sign structure; divergence correction redistributes the
imbalance subject to \(\|\phi\|_\infty \le 1\).

\textbf{Interpretation.} The propagation operator \(d_0 d_0^*\) acts as
a discrete Laplacian, but driven by the \(\ell^1\) subgradient rather
than \(\ell^2\) diffusion. This distinction is fundamental: \(\ell^2\)
diffusion spreads mass uniformly and destroys structure, while
\(\ell^1\) propagation is \emph{edge-selective}, preferring minimal
paths and concentrating signal flow along sharp boundaries.

\textbf{Corollary 4.2 (Finite-Speed Propagation).} A disturbance at
vertex \(x\) affects only vertices within graph distance \(k\) after
\(k\) iterations.

This enforces causal structure in the weak sense: local updates dictate
finite propagation speed, eliminating instantaneous non-local effects.
The iteration index \(k\) serves as the intrinsic time parameter,
corresponding to descent in the \(\ell^1\) action rather than an
externally imposed clock. However, this causal structure does not by
itself imply an invariant propagation speed, Lorentz symmetry, or
observer-independent frame structure; propagation speed depends on the
local graph connectivity and step-size parameters. The framework derives
causal ordering from variational descent; metric spacetime structure
remains an open emergence problem.

\emph{(Cycle--Field Interaction).} When a propagating disturbance
encounters a stable cycle (Theorem 4.5), the cycle resists absorption
due to its divergence-free certification. The disturbance routes
\emph{around} the cycle or modifies only the adjacent lattice. Stable
cycles therefore act as topological obstacles to propagation---the first
structural instance of field--matter interaction derived purely from
\(\ell^1\) constraint logic.

\textbf{Remark 4.3 (Cone Structure from Dual Feasibility).} The dual
feasibility constraint \(\|\phi\|_\infty \le 1\) imposes a bounded flux
per edge per iteration. Combined with locality, this induces a cone of
admissible influence: after \(k\) steps from vertex \(x\), only vertices
within the set \(\mathcal{C}(x,k) = \{ y \mid d_G(x,y) \le k \}\) can be
affected, and the magnitude of influence is bounded by the cumulative
flux capacity along any path. Because \(\ell^1\) dynamics prefer
minimal-support flows (sparse paths), propagation concentrates along
lowest-cost routes rather than spreading isotropically. The defect field
itself naturally defines a path cost
\(d_\delta(x,y) = \inf_{\gamma: x \to y} \sum_{e \in \gamma} |\delta(e)|\),
yielding an emergent metric structure in which geometry depends on the
field configuration. This provides a structural basis for causal
geometry arising from the dynamics, rather than being imposed a priori.
However, this cone and metric are graph-dependent and generically
anisotropic; whether they can yield isotropic large-scale geometry under
aggregation is an open problem (see Section 6.4).

\subsubsection{4.6 The Continuum Limit: Total Variation
Flow}\label{the-continuum-limit-total-variation-flow}

The discrete \(\ell^1\) autopoietic rules provide an exact computable
basis for generating causal interacting poset graphs. We now establish
their formal continuum scaling. The logical chain is: \(\ell^1\) forces
sparsity; sparsity forces concentration of gradients into measure-valued
distributions; measure-valued gradients live in the space of functions
of bounded variation (BV); and the natural variational functional on BV
is total variation. The \(\ell^1\) coboundary structure therefore is
consistent with total variation--type continuum functionals under
standard scaling limits. The precise isotropy of the limiting functional
depends on the symmetry of the underlying discretization: anisotropic or
crystalline graph sequences yield anisotropic TV, while isotropic
refinement yields the standard (Euclidean) total variation. The
continuum geometric structure is inherited from the embedding of the
graph sequence; the \(\ell^1\) framework determines the variational
form, while the ambient geometry determines its metric structure.

\textbf{Theorem 4.8 (Continuum Emergence).} \textbf{(Layer C --- Formal
Limit)} Let \(G_h\) be a sequence of graphs with uniformly bounded
degree admitting a consistent geometric embedding in a domain
\(\Omega \subset \mathbb{R}^d\), with appropriately scaled edge weights
and mesh spacing \(h \to 0\). Under isotropic refinement, the rescaling
\(I_h(s_h) = \sum_e |\delta_h(e)| \to \int_\Omega |\nabla s(x)| \, dx\)
admits total variation flow as the canonical local continuum limit, and
the Local \(\ell^1\) Primal--Dual Iterator is consistent with a
minimizing-movement discretization of this flow:
\[ \partial_t s(x,t) \in -\partial \int_\Omega |\nabla s(x,t)| \, dx \]

Where the field is smooth, this reduces to:
\[ \partial_t s = \nabla \cdot \left( \frac{\nabla s}{|\nabla s|} \right) \]

\textbf{Interpretation.} The iterator defines an intrinsic time
parameter corresponding to descent in the \(\ell^1\) action, yielding a
minimizing-movement interpretation of evolution (in the sense of De
Giorgi), rather than a direct time discretization of the PDE. In the
continuum limit, the \(\ell^1\) defect dynamics are consistent with
total variation flow, producing nonlinear, edge-preserving field
evolution. This establishes a fundamental structural contrast:

\begin{enumerate}
\def\labelenumi{\arabic{enumi}.}
\tightlist
\item
  \textbf{Classical \(\ell^2\) Diffusion:} \(\partial_t s = \Delta s\).
  Isotropically smooths all gradients, destroying object structure
  toward featureless equilibrium.
\item
  \textbf{Autopoietic \(\ell^1\) Flow:}
  \(\partial_t s = \nabla \cdot (\nabla s / |\nabla s|)\). Preserves
  step-edges, concentrates gradients into sharp topological boundaries,
  and drives geometry under mean curvature.
\end{enumerate}

Crucially, the dynamics do not eliminate structure globally; they
concentrate irreducible defect onto stable sets (cycles in the discrete
setting, interfaces in the continuum), preserving topological
information while reducing total variation. Discrete cycles
characterizing autopoietic stable objects converge to sharp level-set
boundaries (codimension-1 manifolds) enclosing persistent domains,
consistent with known graph-to-TV \(\Gamma\)-convergence results in the
variational literature (cf.~Chambolle 2004, van Gennip--Bertozzi 2012).
In the operator-valued case, the Schatten-1 functional
\(\int_\Omega \|\nabla \rho(x)\|_1 \, dx\) suggests a formally analogous
structure, though a rigorous continuum theory for operator-valued total
variation remains open and constitutes a nontrivial extension beyond the
scalar setting.

\subsubsection{\texorpdfstring{4.7 Fiber-Wise Symmetry and \(U(n)\)
Gauge
Fields}{4.7 Fiber-Wise Symmetry and U(n) Gauge Fields}}\label{fiber-wise-symmetry-and-un-gauge-fields}

Because the Local \(\ell^1\) Primal--Dual Iterator continuously strips
away exact artifacts to isolate topological frustration, it enforces
strict constraints on the emergent symmetry groups. We impose a single
mathematically necessary constraint ensuring that optimization logic
remains stable across equivalent observers:

\textbf{Axiom (Local Compositional Invariance).} The absolute \(\ell^1\)
defect magnitude must remain invariant under local relabelings of
internal states. Let \(s(x) \mapsto g_x s(x)\). Preserving the
coboundary norm requires:
\[ \|\delta_e(g \cdot s)\|_1 = \|\delta_e(s)\|_1 \]

This invariance immediately forces the naturality constraint
\(r_{b,a} \circ g_b = g_a \circ r_{b,a}\).

\textbf{Lemma 4.9 (Trace-Norm Automorphisms).} Any surjective trace-norm
preserving, completely positive, compositional automorphism of
\(\mathcal{B}_1(\mathcal{H})\) is unitary conjugation. \emph{Proof
Sketch.} By the Kadison isometry theorem {[}9{]} and its extension to
Schatten ideals by Sourour {[}10{]}, every surjective linear isometry of
\(\mathcal{B}_1(\mathbb{C}^n)\) has the form \(A \mapsto UAV\) or
\(A \mapsto UA^T V\) for unitaries \(U, V\). The requirement of complete
positivity excludes the transpose (which is positive but not completely
positive for \(n \ge 2\), as its Choi matrix has a negative eigenvalue).
The requirement of compositional consistency over the presheaf
restriction maps (\(g_a \circ r_{b,a} = r_{b,a} \circ g_b\)) forces
\(V = U^*\): when \(r_{b,a}\) includes partial traces, the equation
\(U_a \operatorname{Tr}_{\bar{a}}(\rho) V_a = \operatorname{Tr}_{\bar{a}}(U_b \rho V_b)\)
for all \(\rho\) requires \(V_a = U_a^*\) (see, e.g., Watrous, \emph{The
Theory of Quantum Information}, Proposition 2.24). The admissible
transformations reduce to unitary conjugation \(A \mapsto UAU^*\).
\(\square\)

\textbf{Theorem 4.10 (Unitary Gauge Structure).} \textbf{(Layer B)} When
stalks are finite-dimensional operator algebras
\(F(x) = \mathcal{B}_1(\mathbb{C}^n)\) evaluated under the Schatten-1
trace norm (Theorem 3.1), the local transformations satisfying
compositional invariance restrict admissible local automorphisms to
unitary conjugation, corresponding to a \(U(n)\) gauge structure.

\begin{enumerate}
\def\labelenumi{\arabic{enumi}.}
\tightlist
\item
  By Lemma 4.9, the automorphisms \(g_x\) must be trace-norm isometries
  acting via unitary conjugation: \(g_x(A) = U_x A U_x^*\).
\item
  The presheaf restriction maps \(r_{b,a}\) behave as \textbf{discrete
  parallel transport}, defining geometric paths across the tensor
  network.
\item
  The optimal terminal defect \(\delta_{\text{irr}}\) produced by the
  Iterator acts exactly as the unresolvable physical gauge curvature
  (holonomy) constrained within the graph.
\end{enumerate}

The Local \(\ell^1\) Primal--Dual Iterator thus dynamically derives
internal \(U(n)\) gauge fields purely from the requirement to minimize
the Schatten-1 sequence distance between adjoining operator states.

\textbf{Remark (U(n) vs.~PU(n)).} The gauge group acts by conjugation
\(A \mapsto UAU^*\), under which the center \(U(1) \subset U(n)\) acts
trivially (scalar phase cancels). The group that acts faithfully on the
stalks is therefore \(PU(n) = U(n)/U(1)\). We write \(U(n)\) to denote
the group of local frame changes (the assignment \(x \mapsto U_x\)),
which is the standard convention in lattice gauge theory where the
\(U(1)\) phase carries physical meaning as the global gauge choice.
Complete positivity is an additional physical axiom motivated by quantum
mechanics (quantum channels must be CP); it is not forced by the
\(\ell^1\) structure alone.

Establishing subgroups \(SU(3) \times SU(2) \times U(1)\) requires
explicit constraints (rank, phase, symmetry breaking) that descend from
specific bounded dual topologies during combinatorial lattice scaling.

\subsubsection{4.8 Geometry Emergence: Manifolds from
Defects}\label{geometry-emergence-manifolds-from-defects}

The continuum limit reveals a profound geometric consequence of
\(\ell^1\) structure.

\textbf{Theorem 4.11 (Interface Emergence).} \textbf{(Layer C)} In the
continuum limit of the \(\ell^1\) defect functional
\(E(s) = \int_\Omega |\nabla s(x)| \, dx\), under standard BV regularity
assumptions, the support of the singular part of \(\nabla s\)
concentrates on codimension-1 sets \(\Sigma \subset \Omega\), which
evolve under total variation flow by mean curvature.

\emph{Proof Sketch.} Total variation minimization selects functions of
bounded variation whose gradients are Radon measures. By the structure
theorem for BV functions (under appropriate regularity), the singular
part of \(\nabla s\) concentrates on rectifiable sets of codimension 1.
Under TV flow, these sets evolve by mean curvature motion (established
in geometric measure theory by the work of Evans--Spruck and
Chen--Giga--Goto). \(\square\)

\textbf{Interpretation.} The \(\ell^1\) framework produces a
\emph{stratified geometry} in which the field is constant almost
everywhere, with all structure concentrated on evolving codimension-1
manifolds \(\Sigma\). The codimension-1 nature of emergent structures
follows from the measure-minimizing property of total variation, which
penalizes surface measure rather than volume. The emergent interfaces
are not static regularity artifacts but dynamically interacting
structures governed by a variational principle and conservation laws
(Theorems 4.14--4.15), elevating them from mere features of BV
regularity to carriers of physical structure. In this framework,
geometry is encoded variationally through the \(\ell^1\) action rather
than through a pre-specified metric tensor. The discrete-to-continuum
correspondence is direct:

\begin{longtable}[]{@{}
  >{\raggedright\arraybackslash}p{(\linewidth - 2\tabcolsep) * \real{0.5000}}
  >{\raggedright\arraybackslash}p{(\linewidth - 2\tabcolsep) * \real{0.5000}}@{}}
\toprule\noalign{}
\begin{minipage}[b]{\linewidth}\raggedright
Discrete
\end{minipage} & \begin{minipage}[b]{\linewidth}\raggedright
Continuum
\end{minipage} \\
\midrule\noalign{}
\endhead
\bottomrule\noalign{}
\endlastfoot
Stable cycles & Codimension-1 interfaces \(\Sigma\) \\
Cycle stability (Thm 4.5) & Interface persistence under curvature
flow \\
Cycle interaction (Thm 4.6) & Interface merging, splitting,
annihilation \\
Tree instability (Cor 4.1) & Smooth regions carry no structure \\
\end{longtable}

\subsubsection{4.9 Emergent Connection--Curvature--Holonomy
Structure}\label{emergent-connectioncurvatureholonomy-structure}

The dual variables and cycle defects of the \(\ell^1\) framework
naturally organize into a connection-like geometric structure.

\textbf{Definition 4.6 (\(\ell^1\) Holonomy).} For a cycle \(C\) in the
graph, define the \(\ell^1\) holonomy:
\[ \mathrm{Hol}(C) := \sum_{e \in C} \delta_{\mathrm{irr}}(e) \]

Since exact parts cancel around any cycle
(\(\sum_{e \in C} d_0 f = 0\)), holonomy depends only on the cohomology
class \([\delta]\).

\textbf{Definition 4.7 (\(\ell^1\) Connection).} Edge restriction maps
\(r_{b,a}\) define local transport rules. The defect
\(\delta_e = r_{b,a}(s(b)) - s(a)\) measures transport failure. The
connection is the collection of local transport operators; the defect is
its curvature.

\textbf{Definition 4.8 (\(\ell^1\) Curvature).} For a cycle \(C\),
define \(F(C) := \mathrm{Hol}(C)\). Curvature is the failure of
consistent transport around a closed loop---structurally identical to
gauge curvature.

\textbf{Theorem 4.12 (Dual Field as Conserved Flux).} \textbf{(Layer B)}
In the continuum limit, the dual variable \(\phi\) defines a
divergence-free unit flux field:
\[ |\phi(x)| \le 1, \quad \nabla \cdot \phi = 0 \] aligned with defect
gradients and supported on the emergent interfaces \(\Sigma\).

\emph{Proof Sketch.} The dual optimality condition \(d_0^* \phi = 0\)
becomes \(\nabla \cdot \phi = 0\) in the continuum. The constraint
\(\|\phi\|_\infty \le 1\) persists under scaling. Since \(\nabla s\)
concentrates on \(\Sigma\), the alignment condition
\(\phi \in \partial |\nabla s|\) forces \(\phi\) to be supported on
\(\Sigma\) with \(\phi\) parallel to the normal direction. \(\square\)

\textbf{Theorem 4.13 (Emergent Gauge Structure).} \textbf{(Layer B)} In
the \(\ell^1\) autopoietic framework: 1. Edge defects define a
\emph{connection} via local transport inconsistency. 2. Cycle defects
define \emph{holonomy} as irreducible accumulated mismatch. 3. Holonomy
defines \emph{curvature}. 4. Dual variables define \emph{conserved flux
fields} compatible with this structure.

We emphasize that this construction does not assume a Lie group-valued
connection or a principal bundle; rather, the
connection--curvature--holonomy structure arises abstractly from defect
transport on presheaf stalks. The parallel to standard gauge theory is
structural, not definitional. Combined with the \(U(n)\) result (Theorem
4.10), in the operator-valued setting the connection acts on trace-class
stalks with unitary gauge symmetry, and holonomy measures the
accumulated Schatten-1 mismatch around closed loops.

\subsubsection{\texorpdfstring{4.10 The \(\ell^1\) Gauge Action and
Variational
Principle}{4.10 The \textbackslash ell\^{}1 Gauge Action and Variational Principle}}\label{the-ell1-gauge-action-and-variational-principle}

We now promote the defect functional to a gauge-invariant action.

\textbf{Definition 4.9 (\(\ell^1\) Gauge Action).} Define the
\textbf{section-level action}: \[ \mathcal{S}(s) := \|d_0 s\|_1 \]

This is not gauge-invariant: under a gauge shift \(s \mapsto s + f\),
the defect transforms as \(d_0 s \mapsto d_0 s + d_0 f\), and
generically \(\|d_0 s + d_0 f\|_1 \neq \|d_0 s\|_1\). The
\textbf{gauge-invariant action} is the quotient norm:
\[ \bar{\mathcal{S}}([\delta]) := \Phi([\delta]) = \min_f \|\delta - d_0 f\|_1 \]

which depends only on the cohomology class \([\delta] \in H^1\).
Physical configurations minimize \(\bar{\mathcal{S}}\); the
section-level action \(\mathcal{S}(s)\) is an upper bound for
\(\bar{\mathcal{S}}([d_0 s])\).

\textbf{Theorem 4.14 (\(\ell^1\) Variational Principle).} \textbf{(Layer
B)} Physical configurations are minimizers of the gauge-invariant
action:
\[ \delta_{\mathrm{irr}} = \arg\min_{[\delta]} \bar{\mathcal{S}}([\delta]) \]

The Euler--Lagrange condition is:
\[ 0 \in \partial \|d_0 s\|_1 \quad \Longleftrightarrow \quad d_0^* \phi = 0, \quad \phi \in \partial \|\delta\|_1 \]

That is: the field equation \emph{is} the divergence-free dual
condition. Conservation laws are field equations.

\textbf{Curvature Formulation.} When the irreducible defect
\(\delta_{\mathrm{irr}}\) is supported on edge-disjoint cycles
\(\{C_i\}\), the gauge-invariant action rewrites as:
\[ \bar{\mathcal{S}} = \sum_i |\mathrm{Hol}(C_i)| \]

For general defect configurations on overlapping cycles, this
decomposition is a lower bound; the exact action requires solving the
\(\ell^1\) quotient norm optimization. The cycle-basis representation is
the structural \(\ell^1\) analogue of the Yang--Mills action
\(\sum \|F\|^2\), with \(\ell^1\) replacing \(\ell^2\).

\textbf{Dual Formulation.} By LP duality:
\[ \mathcal{S}(\delta) = \max_{\phi:\, d_0^* \phi = 0,\, \|\phi\|_\infty \le 1} \langle \phi, \delta \rangle \]

The action equals the maximum flux compatible with divergence-free and
boundedness constraints. This primal--dual identity---action equals flux
capacity---is the structural core of the framework.

\textbf{Continuum Action.} Under the scaling of Theorem 4.8:
\[ \mathcal{S}(s) = \int_\Omega |\nabla s| \, dx \] whose variation
recovers TV flow, confirming consistency with the iterator dynamics.

\subsubsection{\texorpdfstring{4.11 Symmetry and Conservation:
\(\ell^1\) Noether
Correspondence}{4.11 Symmetry and Conservation: \textbackslash ell\^{}1 Noether Correspondence}}\label{symmetry-and-conservation-ell1-noether-correspondence}

\textbf{Definition 4.10 (Symmetry of the \(\ell^1\) Action).} A
transformation \(T\) acting on sections \(s \in C^0\) is a
\emph{symmetry} if \(\|d_0(Ts)\|_1 = \|d_0 s\|_1\). Examples include
additive gauge shifts (\(s \mapsto s + f\) for \(f\) constant on
connected components), unitary conjugation (operator-valued case,
Theorem 4.10), and graph automorphisms.

\textbf{Theorem 4.15 (\(\ell^1\) Noether-Type Conservation Law).}
\textbf{(Layer B)} Let \(T_\epsilon\) be a continuous one-parameter
family of symmetries of the \(\ell^1\) action. Then along any solution
of the \(\ell^1\) variational principle, there exists a dual field
\(\phi\) such that:
\[ d_0^* \phi = 0 \quad \text{and} \quad \left.\frac{d}{d\epsilon}\right|_{\epsilon=0} \|d_0(T_\epsilon s)\|_1 = \langle \phi, d_0(\dot{s}) \rangle = 0 \]

\emph{Proof Sketch.} Action invariance gives
\(\frac{d}{d\epsilon} \mathcal{S}(T_\epsilon s) = 0\). By the
subgradient structure,
\(\delta \mathcal{S} = \langle \phi, \delta(\cdot) \rangle\) for
\(\phi \in \partial \|\delta\|_1\). The stationarity condition
\(\langle \phi, d_0 \dot{s} \rangle = 0\) integrates by parts (discrete
adjoint) to \(\langle d_0^* \phi, \dot{s} \rangle = 0\). Since this
holds for all admissible variations \(\dot{s}\), we obtain
\(d_0^* \phi = 0\). \textbf{Regularity caveat:} The \(\ell^1\) norm is
not differentiable at points where \(\delta_e = 0\) for some edges. At
such points, the subdifferential \(\partial \|\delta\|_1\) is
multi-valued (the dual variable \(\phi(e) \in [-1, 1]\) is not uniquely
determined). The conservation law holds for every subgradient selection
\(\phi \in \partial \|\delta\|_1\), but the conserved dual field is not
unique at non-smooth configurations. \(\square\)

\textbf{Interpretation.} Symmetries of the \(\ell^1\) action generate
divergence-free dual fields---conserved flux. In classical physics,
Noether's theorem says: symmetry \(\to\) conserved current. In the
\(\ell^1\) framework: symmetry \(\to\) dual feasibility \(\to\)
conserved flux. The correspondence is structural, not analogical.

\begin{longtable}[]{@{}ll@{}}
\toprule\noalign{}
Classical Field Theory & \(\ell^1\) Cohomological Framework \\
\midrule\noalign{}
\endhead
\bottomrule\noalign{}
\endlastfoot
Action \(\mathcal{S}\) & \(\|d_0 s\|_1\) \\
Field & Section \(s\) \\
Curvature & Holonomy \(\mathrm{Hol}(C)\) \\
Gauge invariance & \(s \sim s + f\) (quotient structure) \\
Field equation & \(d_0^* \phi = 0\) \\
Conserved current & Divergence-free \(\phi\) \\
\end{longtable}

\textbf{Continuum Version.} Under scaling, the conservation law becomes
\(\nabla \cdot \phi = 0\) with \(\phi\) the continuum dual flux. This is
the classical conservation equation structure.

\begin{center}\rule{0.5\linewidth}{0.5pt}\end{center}

\subsection{5. The Distinguishability Bound (The Bridge to
Physics)}\label{the-distinguishability-bound-the-bridge-to-physics}

\subsubsection{5.1 The Distinguishability Bound
Theorem}\label{the-distinguishability-bound-theorem}

The results of Papers 000--004 establish that the \(\ell^1\) coboundary
norm is the unique additive defect diagnostic, and that it lifts
canonically to the Schatten-1 trace norm natively in operator-valued
presheaves. This structural rigidity allows us to formally map abstract
topological obstruction directly to standard distance metrics in
information theory.

\begin{quote}
\textbf{Theorem 5.1 (Distinguishability Bound).} \textbf{(Layer B)} Let
\((\rho_x)_{x \in P}\) be an operator-valued presheaf section over a
path graph connecting nodes \(a\) and \(b\), with uniform local Hilbert
spaces \(\mathcal{H}_x = \mathcal{H}\) for all \(x\), and let
\(I(s) = \sum_e \|\delta_e(s)\|_1\) be the \(\ell^1\) coboundary defect.
Let \(R_{b,a} = r_{1,0} \circ r_{2,1} \circ \cdots \circ r_{m,m-1}\)
denote the composed restriction (parallel transport) from \(b\) to
\(a\). Then the quantum trace distance satisfies:
\[ T(\rho_a, R_{b,a}(\rho_b)) \le \frac{1}{2} I(s) \]
\end{quote}

\emph{Proof Sketch.} By telescoping and the contractivity of each
restriction map:
\(\|R_{b,a}(\rho_b) - \rho_a\|_1 \le \sum_{i} \|r_{i+1,i}(\rho_{i+1}) - \rho_i\|_1 = \sum_{e \in \text{path}(a,b)} \|\delta_e\|_1 = I(s)\).
Noting that trace distance is precisely half the trace norm yields the
result. When all Hilbert spaces are identical and restriction maps are
the identity, this reduces to the standard triangle inequality
\(T(\rho_a, \rho_b) \le \frac{1}{2} I(s)\). \(\square\)

This theorem provides a structural link between the framework and
operational distinguishability in quantum information, contingent on
identifying presheaf sections with physical quantum states. The
\(\ell^1\) topological tension \(I(s)\) actively bounds the maximum
quantum distinguishability between any two local subsystems.
Furthermore, returning to the Hodge cohomology of Paper 003, the
quotient norm \(\Phi([\delta])\) bounds the \emph{minimal unavoidable
distinguishability} (irreducible correlation) that remains after all
local gauge transformations are exhausted.

\subsubsection{5.2 Connection to the Fermion Sign Problem (Layer C ---
Conjectural)}\label{connection-to-the-fermion-sign-problem-layer-c-conjectural}

The Schatten-1 extension (Theorem 3.1) suggests a structural connection
to the fermion sign problem. When the presheaf stalks carry fermionic
statistics, the restriction maps \(r_{b,a}\) involve antisymmetrization,
and the coboundary defect \(\delta_e\) would measure the failure of
local fermionic states to glue consistently. Whether the trace norm of
this defect can be rigorously identified with sign frustration requires
formal development beyond this paper's scope {[}11{]}.

If this connection holds, it would provide a structural interpretation:
the fermion sign problem would be a manifestation of \(\ell^1\) gauge
frustration on operator-valued presheaves, rather than an artifact of
computational representation. However, this remains conjectural and
requires explicit construction of the fermionic presheaf and
verification that the coboundary defect captures the relevant sign
structure.

\subsubsection{5.3 Connection to Coupling Constraints (Layer C ---
Conjectural)}\label{connection-to-coupling-constraints-layer-c-conjectural}

Subsequent work {[}12{]} explores whether the fine structure constant
can be derived from the bounding limits of a minimal \(\ell^1\)
obstruction seed. The structural constraints proposed here suggest that:

\begin{enumerate}
\def\labelenumi{\arabic{enumi}.}
\tightlist
\item
  The \(\ell^1\) obstruction on a 2-node graph has an absolute
  coboundary norm gap.
\item
  Subdivision interpolation (expanding to a 3-node lattice) preserves
  the defect monotonically by Theorem 2.1.
\item
  The Schatten-1 extension (Theorem 3.1) guarantees this topological
  constraint survives when evaluating localized quantum states.
\end{enumerate}

Whether the derivation of wave mechanics, the Born rule, and coupling
constants formally descends from these bounding trace constraints
remains deeply speculative. The present paper only provides the
structural foundation demonstrating that the \(\ell^1\)
Distinguishability Bound mathematically governs such limiting
constraints.

\begin{center}\rule{0.5\linewidth}{0.5pt}\end{center}

\subsection{6. Discussion}\label{discussion}

\subsubsection{6.1 What This Paper Proves}\label{what-this-paper-proves}

This paper extends the foundational quad (Papers 000--003) in five
directions:

\begin{enumerate}
\def\labelenumi{\arabic{enumi}.}
\item
  \textbf{Functorial aggregation is monotone} (Theorem 2.1): \(\ell^1\)
  defects cannot be hidden by coarse-graining.
\item
  \textbf{The operator-valued extension is canonical} (Theorem 3.1): The
  trace norm is the unique diagnostic on operator-valued presheaves
  satisfying the foundational axioms.
\item
  \textbf{The Local \(\ell^1\) Primal--Dual Iterator characterizes
  defect fixed points} (Definition 4.1, Theorem 4.2, Theorem 4.4): A
  local \(\ell^1\) primal--dual iterator whose fixed points are
  \(\ell^1\)-minimal representatives (Theorem 4.2), with edge-wise local
  descent bounds (Theorem 4.4). Monotonic global descent is conjectured
  (Conjecture 4.3) but not proved.
\item
  \textbf{Stable structure emerges and interacts} (Theorems 4.5--4.7):
  Cycle-supported defects are fixed points; overlapping cycles interact
  via re-optimization; disturbances propagate at finite speed.
\item
  \textbf{A gauge-invariant variational structure exists} (Theorems
  4.11--4.15): The \(\ell^1\) action, its curvature formulation, the
  emergent connection--holonomy framework, and the Noether-type
  conservation law provide a field-theoretic structure derived from
  defect minimization alone.
\end{enumerate}

\subsubsection{6.2 What This Paper Does Not
Prove}\label{what-this-paper-does-not-prove}

The conjectural elements (Layer C) are explicitly marked:

\begin{itemize}
\tightlist
\item
  The continuum limit (Theorem 4.8) is stated as a formal scaling limit;
  full \(\Gamma\)-convergence is not proved.
\item
  The geometry emergence result (Theorem 4.11) invokes established
  results from geometric measure theory under standard BV regularity
  assumptions but does not re-derive them.
\item
  The connections to the Standard Model gauge groups remain structural
  observations pending explicit derivation.
\item
  The derivation chain to quantum mechanics (Section 5.1) contains both
  established and open steps.
\item
  Monotonic descent of the iterator (Conjecture 4.3) is not proved; the
  convergence analysis for the three-step primal--dual variant is
  deferred to subsequent work.
\item
  The finite-speed propagation (Corollary 4.2) establishes causal
  ordering but does not imply Lorentz invariance, an invariant
  propagation speed, or relativistic frame equivalence; whether
  additional symmetry constraints on the \(\ell^1\) action can force an
  invariant speed remains open.
\end{itemize}

The conditionality of all claims is preserved. The \(\ell^1\) framework
is forced \emph{given} the axioms of Papers 001--002. The extensions
proved here are forced \emph{given} the \(\ell^1\) framework. The
conjectural elements indicate where the program requires further formal
development.

\subsubsection{6.3 Position in the
Program}\label{position-in-the-program}

Papers 000--003 answer: \emph{What must the defect metric be?}
(\(\ell^1\), uniquely.)

This paper answers: \emph{What is its static and variational structure?}
(Monotone under aggregation, canonical on operators, governed by a
gauge-invariant action whose symmetries generate conserved flux, with
emergent connection--curvature--holonomy structure and a formal
continuum limit recovering total variation flow.)

Paper 005 inherits: \emph{How does local dynamics realize this action?}
(The Autopoietic Iterator as the local engine driving causal poset
construction, propagation, and structural emergence.)

The program is falsifiable at every joint: if any proved theorem is
incorrect, or if any conjectured step fails, the downstream structure
requires revision at that point. The separation of established results
(Layers A, B) from conjectures (Layer C) is maintained to make this
falsification structure transparent.

\subsubsection{6.4 The Scope of Geometric
Emergence}\label{the-scope-of-geometric-emergence}

The \(\ell^1\) framework does not prescribe the ambient geometric
substrate---the topology, dimensionality, or embedding of the underlying
graph or poset are inputs, not outputs. Rather, the framework determines
the variational structure governing defect minimization, which in turn
forces all nontrivial structure to concentrate on minimal geometric
carriers. In the continuum limit, this yields codimension-1 interfaces
evolving under curvature-driven dynamics. Geometry in this sense is not
imposed a priori but emerges as the support of irreducible defect under
the \(\ell^1\) action.

Concretely, the framework derives:

\begin{itemize}
\tightlist
\item
  \textbf{Where structure lives:} on minimal-measure carriers (cycles →
  interfaces)
\item
  \textbf{How structure behaves:} curvature-driven evolution,
  finite-speed propagation
\item
  \textbf{What counts as irreducible:} gauge-invariant cohomological
  content (\(\Phi([\delta])\))
\end{itemize}

It does not derive:

\begin{itemize}
\tightlist
\item
  The ambient dimension or topology of the substrate
\item
  A Riemannian metric on the background space
\item
  The specific graph connectivity or embedding
\end{itemize}

This distinction positions the framework as a theory of
\textbf{geometric content} (what structures exist and how they interact)
rather than \textbf{geometric background} (what space they inhabit). The
substrate is given; the physics living on it is derived. The framework
derives causal structure first (finite-speed propagation, causal
ordering from variational descent); metric structure---including Lorentz
symmetry, invariant speed, and observer equivalence---remains an open
emergence problem requiring additional symmetry constraints on the
action or the graph ensemble.

\subsubsection{6.5 Open Problems: Emergent Metric and
Homogenization}\label{open-problems-emergent-metric-and-homogenization}

The cone structure identified in Remark 4.3 and the defect-induced path
cost \(d_\delta(x,y) = \inf_\gamma \sum_{e \in \gamma} |\delta(e)|\)
define a candidate emergent metric. Three open questions determine
whether this candidate can bridge to physical geometry:

\begin{enumerate}
\def\labelenumi{\arabic{enumi}.}
\item
  \textbf{Homogenization under aggregation.} While the defect-induced
  metric is generically anisotropic at microscopic scales, the
  combination of \(\ell^1\) minimization (which penalizes redundant
  directional concentration) and functorial aggregation (Theorem 2.1,
  which preserves total defect under coarse-graining) suggests a
  homogenization mechanism. \emph{Open problem:} Does repeated
  coarse-graining of generic \(\ell^1\) defect configurations converge
  to an approximately isotropic metric in distribution, for sufficiently
  symmetric or ergodic graph ensembles?
\item
  \textbf{Invariant speed.} The dual feasibility constraint
  \(\|\phi\|_\infty \le 1\) imposes a uniform flux bound, yielding a
  maximal propagation rate. \emph{Open problem:} Under what symmetry
  conditions on the \(\ell^1\) action does this maximal rate become
  invariant under admissible transformations?
\item
  \textbf{Metric--dynamics coupling.} The defect-induced metric
  \(d_\delta\) depends on the field configuration, so geometry and
  dynamics are coupled. \emph{Open problem:} Does this coupling admit a
  consistent continuum limit in which the emergent metric satisfies a
  geometric evolution equation (analogous to Ricci flow or Einstein
  equations)?
\end{enumerate}

A positive answer to (1) would establish a serious bridge between the
\(\ell^1\) framework and spacetime structure. A counterexample---a
stable, strongly anisotropic defect configuration resistant to
aggregation---would bound the geometric ambitions of the program. This
constitutes the critical falsification test for the physics-level
extension of the framework.

\begin{center}\rule{0.5\linewidth}{0.5pt}\end{center}

\subsection{Acknowledgments}\label{acknowledgments}

No external funding. No conflicts of interest.

\begin{center}\rule{0.5\linewidth}{0.5pt}\end{center}

\subsection{References}\label{references}

{[}1{]} S. Mac Lane, \emph{Categories for the Working Mathematician},
Springer, 1971.

{[}2{]} B. Eckmann, ``Harmonische Funktionen und Randwertaufgaben in
einem Komplex,'' \emph{Commentarii Mathematici Helvetici} \textbf{17}
(1944/45), 240--255.

{[}3{]} J. A. Hansen and R. Ghrist, ``Toward a spectral theory of
cellular sheaves,'' \emph{Journal of Applied and Computational Topology}
\textbf{3} (2019), 315--358.

{[}4{]} J. H. Carroll, ``Projection Obstruction Theory - Retraction
Nonlinearity, l1 Rigidity, and Density Scaling,'' Zenodo, 2026.
https://doi.org/10.5281/zenodo.19151803

{[}5{]} J. H. Carroll, ``The Free \(\ell^1\) Seminorm on Banach Presheaf
Coboundaries,'' Zenodo, 2026. https://doi.org/10.5281/zenodo.19152115

{[}6{]} J. H. Carroll, ``Coordinate-Wise Additivity and the \(\ell^1\)
Norm on Finite Graph Cochains,'' Zenodo, 2026.
https://doi.org/10.5281/zenodo.19152599

{[}7{]} J. H. Carroll, ``Hodge Structure and Gauge Equivalence of
\(\ell^1\) Defect Fields,'' Zenodo, 2026.
https://doi.org/10.5281/zenodo.19152799

{[}8{]} R. T. Rockafellar, \emph{Convex Analysis}, Princeton University
Press, 1970.

{[}9{]} R. V. Kadison, ``Isometries of operator algebras,'' \emph{Annals
of Mathematics} \textbf{54} (1951), 325--338.

{[}10{]} A. R. Sourour, ``Isometries of norm ideals of compact
operators,'' \emph{Journal of Functional Analysis} \textbf{43} (1981),
69--77.

{[}11{]} M. Troyer and U.-J. Wiese, ``Computational complexity and
fundamental limitations to fermionic quantum Monte Carlo simulations,''
\emph{Physical Review Letters} \textbf{94} (2005), 170201.

{[}12{]} J. H. Carroll, ``Emergent Gauge Symmetry from \(\ell^1\) Defect
Flow,'' Zenodo, 2026. https://doi.org/10.5281/zenodo.19182289

\end{document}
